\begin{studiumworked}{M09. Divergencia y definición de rendimiento · Intermedio}
\textbf{Enunciado.} Un modelo de descarga cónica uniforme por ángulo sólido tiene semiángulo de 14 grados. El empuje cinético de referencia es 3000 N y el término de presión 120 N. Halle el empuje axial corregido. Si un coeficiente de velocidad independiente vale 0.96, determine su rendimiento cinético equivalente, sin aplicarlo otra vez al mismo resultado.

\textbf{Datos.} Semiángulo: $14^\circ$; Momento de referencia: $3000\,\mathrm{N}$; Presión: $120\,\mathrm{N}$.

\textbf{Incógnitas.} Magnitudes solicitadas y cierre físico.

\textbf{Hipótesis.} Modelo académico con hipótesis del capítulo. Datos originales; no representan hardware ni ensayo.

\textbf{Modelo físico.} Conservación y relaciones constitutivas declaradas.
\[\lambda=(1+\cos\alpha)/2,\quad F=\lambda F_{cin}+F_p\]
\[\eta_n=C_v^2\]
\textbf{Desarrollo.}
\begin{enumerate}
\item Integrar $\cos\theta$ con peso $\sin\theta$ da $\lambda=0.9851479$. Se corrige sólo momento axial: $0.9851479(3000)+120=3075.444$ N.
\item La comparación de velocidades independiente tiene $\eta_n=0.96^2=0.9216$. No es 0.96 si el rendimiento se define por energía cinética.
\item No se multiplican ambos factores sobre un empuje ya medido o corregido sin identificar qué pérdidas contiene cada referencia.
\end{enumerate}
\textbf{Operaciones para contrastar el cálculo.}
\begin{enumerate}

\item Magnitud 1.

\[\fitmath{y_{1}=\frac{1+\cos\left(\frac{14\cdot \pi}{180}\right)}{2}\cdot 3000+120\approx 3075.444\,\mathrm{N}}\]

\item Magnitud 2.

\[\fitmath{y_{2}={0.96}^{2}\approx 0.9216\,}\]

\end{enumerate}

\textbf{Resultados.} $F=3075.444$ N; $\eta_n=0.9216$.

\textbf{Interpretación.} El resultado corresponde a las hipótesis indicadas; su interpretación requiere conservar fronteras y unidades.

\textbf{Comprobación.} Para semiángulo cero, lambda vale uno y se recuperan 3120 N. El rendimiento energético permanece menor que el coeficiente de velocidad para valores entre cero y uno.

\textbf{Límites.} El reparto angular es una idealización; no se predice una tobera real de campana.
\end{studiumworked}
